\section{Conclusion}
\label{sec:conclusion}

This paper asked four questions of a single fixed panel of 218{,}934
thirty-minute S\&P~500 realized-variance bars. It answers a fifth it did not
ask, and that answer subsumes much of the rest: almost all of the modelling
apparatus we built can be deleted without loss, and the effects that remain are
smaller than the protocol and data-quality effects surrounding them.

\emph{What should the instrument be?} A truncated power-law kernel with a
clock organ. The fitted HAR coefficients imply a per-lag effective kernel that
is approximately power-law on log-log axes, which recasts the lag ladder as a
quadrature grid rather than a set of economically distinct horizons. The
recasting is not decorative: it predicts that refining the grid pays, and the
grid refines monotonically across six geometric bases, with base-2 beating the
base-5 incumbent by $0.00082$ at DM $t = -10.6$ and the conventional
daily/weekly/monthly structure faring worse than either. It also predicts that
a power law should beat an exponential, which it does by $0.011$ in forecast
loss and decisively in kernel shape. Because the shared benchmark uses the
coarser grid, we carry an explicit $0.00082$ of benchmark debt through every
downstream increment rather than absorbing it.

The self-exciting reading that the power-law shape invites is tested rather
than asserted, and it splits in two. The shape survives: the decay exponent,
estimated across 41 exogenous channels where the log-log slope is invariant to
channel scaling, is $1.033$ with a channel-clustered standard error of $0.036$,
agreeing with the Hawkes-in-finance literature and rejecting both the $1.36$
earlier drafts had carried as a plotting constant and the $\approx 1.6$ that a
rough-volatility correspondence implies. The sign structure appeared to fail, and
explaining why produced the section's most transferable result. The estimate
this project had been reading conditions the self-kernel on 41 exogenous
channels, several of which are direct volatility measures; those absorb the
target's persistence and leave an integrated kernel of $-0.363$ with 14\%
negative mass. The marginal self-kernel on identical rows has integrated mass
$0.913$, no negative mass, and monotone power-law decay at
$\hat\beta = 1.506 \pm 0.035$ over ten lags---sub-critical, close to unity, and
squarely where the reflexivity literature places financial activity. A
self-kernel estimated alongside volatility proxies is not a self-kernel, and
the fit gives no signal that anything is wrong.

\emph{Where does the information live?} Everywhere and nowhere in particular.
Every exogenous bucket improves on the backbone with overwhelming individual
significance, yet the buckets overlap so heavily that the sum of within-class
gains is $2.08$ times the joint gain. Within the least-squares class, 96\% of
explainable variance belongs to the backbone; the 359-column exogenous block
adds a unique contribution that a single mechanism-targeted engineered feature
exceeds by roughly two orders of magnitude per column---a ratio we report with
its arithmetic and its caveats rather than as an effect size. The individual
feature landscape repeats the pattern one level down: one dominant cheap
feature, one partially redundant companion, a long tail of tiny effects, and a
genuine null.

\emph{Does extraction require nonlinearity?} Yes, but boundedly, and the bound
is the interesting part. Tree ensembles beat the best penalized linear model on
all nine feature sets, by between $0.00184$ and $0.00358$; the consistent sign
is the evidence, and the pairwise difference remains untested. The pure
nonlinearity premium on backbone features is $0.00215$ measured against a
correctly specified linear ladder, some 28\% below the figure obtained against
the incumbent. Tuning is a non-lever from every angle: causally tuned linear
arms are worse than the untuned incumbent, and a test-set oracle over penalties
moves it by at most $0.00015$. What the trees exploit saturates at
additive-plus-pairwise structure anchored to the session clock, and deeper
capacity, deep sequence models, and mixture-of-experts routing all fail to
improve on it---a ceiling that a quadratic self-exciting feedback predicts. The
resulting ordering---information $>$ model class $>$ hyperparameters, by
roughly a factor of three at each step---is the paper's central practical law.

\emph{What algorithm should do the extracting?} With features frozen,
exactness and anchoring matter more than expressiveness. Rank-1 rolling least
squares makes the every-bar OLS benchmark computable at scale; the online
elastic-net homotopy extends the same exactness to $\ell_1$ estimation; and the
causal tuning protocol converts ``tuning does not help'' from an anecdote about
one search into a claim with a defensible causal structure. The spectral
$k$-nearest-neighbors proposal is instructive precisely because its ablations
are decisive against it: the OLS anchor is load-bearing (worth up to $0.0132$),
large-neighborhood averaging is load-bearing (worth $\approx 0.0018$), and the
manifold embedding---the novel part---hurts on all eight matched
configurations. A matched retrieval-metric experiment then locates the cause,
and it is not the out-of-sample extension error we had blamed: a
principal-component projection at the same dimension has no extension step and
loses nearly as much, while a supervised coordinate of that dimension ties the
full ambient view. What fails in this regime is unsupervised geometry, at any
dimension and whether linear or nonlinear---and, relatedly, breadth in a
distance, since dropping the exogenous block entirely improves retrieval by
$0.01476$ ($t = -15.5$). A fitted map with learned weights profits from many
weak channels; a metric that weights every coordinate equally is buried by
them. What survives is an unembedded, anchored analog search that beats the
incumbent at QLIKE $0.13096$; it also beats 25 of the 45 causally tuned
parametric arms, while trailing the best arms on the informative buckets. In
the dense-but-weak regime, nonparametric flexibility is a cost to be justified,
not a default.

\emph{How much of the apparatus is load-bearing?} Almost none of it. The
$41 \times 12$ exogenous operator, its singular value decomposition, the
session-conditional frames, the CP and Tucker factorisations, the
rank-constrained priors and the bespoke projection penalty are together
reproduced by contracting each channel's rung ladder against a small number of
fixed shapes and running plain ridge on the result. On the uncorrected panel
that compression appeared to \emph{beat} the full design---$53$ columns at
$0.13280$ against $0.13617$ for plain ridge on all $504$ features, itself worse
than the $12$-column backbone at $0.13571$---and we built a draft around it. On
the corrected panel the full design reaches $0.13146$ and wins, our
pre-registered falsification condition for the compressed one is met at
$-0.00126$ ($t = -5.91$), and the compression is withdrawn
(Section~\ref{sec:ledger}).

What replaces it is a tie. Sweeping rank on the lag axis under two independent
orthonormal bases, the curve is flat from $r = 4$ upward, and the best
compressed arm---$135$ columns---reaches $0.13117$ against $0.13146$ for all
$504$. That $0.00029$ does not survive a widened HAC bandwidth, a block
bootstrap, an era-level test or a leave-one-era-out sweep, so we report the two
as equivalent. A nearly fourfold difference in design width buys nothing
measurable, which is this paper's thesis applied to its own modelling section.
The effective-degrees-of-freedom accounting says why: the elaborate penalty
spends $\approx 54$ degrees of freedom, which is $1 + 12 + 41$, the dimension of
its own null space. The compressed design is an order of
magnitude more robust in the one era where the exogenous block fails
($+0.00347$ against $+0.03321$). The generalisation we are willing to defend is
that in this regime every \emph{estimated} structure is a cost requiring
evidence: eigenmaps, principal components, the operator's own leading direction
re-estimated or frozen, and rank-one factors with fitted loadings all lost to
fixed alternatives, without exception.

\emph{What dominates all of it?} Protocol and data quality---where they bind,
by orders of magnitude; on average, by less than we first claimed. A fixed
training window rather than a rolling one is worth $0.11$ QLIKE. The imputation
defect of Section~\ref{sec:stale_data}---an availability indicator computed
after a forward fill, which silently disabled three guards this codebase
already contained for that exact pathology---put scaled features at $2{,}260$
rolling interquartile ranges and cost a factor of fifty on its window,
$6.37170$ against $0.12391$. Selecting a shape parameter against the reported
loss rather than a validation tail is worth $0.0033$. Against these, the
model-class effects this literature argues about are $0.002$ to $0.004$.

We pre-registered the stronger form of this claim and it failed. Pooled over
six designs, the data-version main effect is $0.00067$ against a design main
effect of $0.00311$: averaged across twenty years and across designs that
mostly do not touch the affected columns, data representation does \emph{not}
dominate. What dominates is the conditional statement---where a fill defect
binds, it is worth more than every modelling choice combined, and nothing in
the fit tells you whether it binds. That is the transferable finding, and it is
narrower and more useful than the one we set out to make: a team that checks
its availability bookkeeping against its raw feeds will not lose fifty-fold to
a dead channel, and no amount of care about model class substitutes for that
check.

\emph{And is any of it stable?} Not as reported. The exogenous block beats the
backbone in seven consecutive eras and fails once, so a twenty-year aggregate
supports opposite conclusions depending on where the window ends, and 27 bars
out of $27{,}376$ carry $47\%$ of the failure. We regard single full-sample
aggregates over multi-decade panels as insufficient reporting, in this paper
and in the literature it follows.

\paragraph{On the evaluation itself.} Two findings concern the measuring
apparatus rather than the models, and we think they travel beyond this panel.
QLIKE and raw-variance MSE rank our 45-arm battery almost independently (rank
correlation $+0.15$, $p = 0.32$, where agreement would give a negative value),
and 14 of the 42 arms that improve on the incumbent under QLIKE are worse than
it under MSE. Separately, QLIKE systematically rewards attenuation: the rank
correlation between QLIKE and the Mincer--Zarnowitz slope is $+0.74$
($p = 5.8 \times 10^{-9}$), so the best-scoring arms are the most
miscalibrated, with the second-best arm in the battery carrying its worst
calibration. Neither observation invalidates a strictly consistent scoring
rule, but both change what a QLIKE ranking should be taken to mean.

\paragraph{On the cost of testing what we claimed.} Section~\ref{sec:inference}
subjects the paper's inherited claims to per-bar inference for the first time.
Four survive Holm-adjusted testing, two are refined, and three do not
survive---trees as residual readers, Fourier clock modulation over session
gates, and rank versus divide were all directionally correct point estimates
that had been promoted to results. A fourth, a singleton model confidence set
over feature sets, turned out never to have been computed at all. We report the
casualty list because it may be the most transferable finding here: on a
2{,}189-bar tile a model confidence set retains 123 of 161 arms, so a tile
leaderboard is not a ranking, and a difference read off one is real only if it
replicates at scale.

\paragraph{What replicated.} The strongest evidence for the paper's central
claims is no longer any single table but Section~\ref{sec:replication}: on a
second basis, with different machinery, under pre-registration with recorded
expectations, the same laws re-derive --- dense-generous-shrunk beats every
selective alternative it met (sixteen more retired arms), the exploitable
content is fast in horizon and rewards per-bar freshness in the first moment
while the second moment obeys neither law, and every estimated geometry
without supervision lost to a frozen one, including the retrieval corrections
whose panel-specific success Section~\ref{sec:knn_selection} could not
confirm. The replication also sharpened the era story: regime change arrived
there not as decay of predictability but as collapse of dictionary
complexity, a form of nonstationarity that headline accuracy metrics are
structurally unable to see.

\paragraph{Open items.} We flag the incomplete edges of the evidence
explicitly, and Section~\ref{sec:limitations} treats them at length.
(i)~The bucket battery has not been re-run on the base-2 backbone; until it is,
every increment carries $0.00082$ of benchmark debt and the smallest effects
are not claimed. (ii)~The model-class gap is now tested pairwise, by a bound that assumes
nothing about the dependence between arms and resolves all nine matched
comparisons after Holm adjustment; the headline pair is the weakest at
$|t| \ge 2.19$, and a model confidence set over the 45 arms remains unavailable
because their per-bar losses were not retained. (iii)~The kernel estimate is reproducible and, once marginal rather than
partial, admissible: $\hat n = 0.913$ with no negative mass. We report it as an
endogeneity analogue rather than an identified branching ratio, because it is
ridge-penalised on transformed data and because a 30-minute aggregated variance
is not the event stream a point-process model describes. Two further
limitations were settled along the way: the 5/95 target winsorization
costs roughly $0.024$ QLIKE and should be expected to disappear from the
pipeline, and the rolling diurnal divisor used alone reaches QLIKE $0.558$
against $1.320$ for an unconditional mean, so it closes about 71\% of the
distance to a fitted model before any model is fitted. (iv)~A
fractional-kernel arm sits in the same power-law family as our backbone and is
absent from the battery; it is the most defensible missing benchmark.
(v)~The winning spectral-$k$NN arm is the minimum of 32 cells selected on the
evaluation panel, and survives on one bucket while sibling buckets failed as
singular; the failure needs diagnosis before the bucket ranking can be trusted,
and the replication panel's pre-registered three-view test declined to confirm
the family (Section~\ref{sec:knn_selection}).
(vi)~The attribution results point outward: the highest-value next investment
is not a larger model but mechanism-targeted data---auction imbalance, order
flow, and dealer positioning aligned to the session-transition regime where the
nonlinear signal concentrates.
